
This paper begins with a simple observation: the aggregation step that converts per-token log-probabilities into a single hallucination confidence score is a universal but unstudied design choice. The central finding is that this choice is not arbitrary --- it is determined by the structural nature of the hallucination being detected, and the wrong choice costs up to 12 AUROC points.

The optimal aggregation function is determined by hallucination type. Recall-failure hallucinations concentrate uncertainty at a single worst-case token; minimum log-probability is maximally discriminative. Imitative-falsehood hallucinations distribute uniform confidence throughout; mean log-probability integrates the diffuse signal. This three-step mechanism is confirmed at distributional, rank-order, and AUROC levels across two model families.

Principal contributions: (1) the first principled aggregation selection criterion for zero-cost hallucination detection, confirmed with large effect sizes on LLaMA-2-7B and Mistral-7B-v0.1; (2) direct distributional evidence for the mechanism's first step (peakedness ratio p = 0.002 on LLaMA-2-7B TriviaQA); and (3) the finding that raw unnormalized log-probability sum is the strongest single-pass factual-recall detector, surpassing published multi-sample Semantic Entropy AUROC at one-tenth the inference cost under this evaluation protocol.

Three immediate extensions are grounded in existing code: NQ completion (inference pipeline complete; GPU time only), length-stratified AUROC analysis (\texttt{length\textbackslash \_stratified\textbackslash \_auroc} implemented in h-e1), and peakedness measurement on TruthfulQA to directly verify the flat-distribution assumption.

The token log-probability sequence contains more structure than is typically exploited. Reading that structure correctly --- by matching the aggregation function to the hallucination type --- costs nothing and yields up to 12 AUROC points.

